\section{Who Can Join, and Who Decides}

\subsection{Eligibility as a two-way gate}

\draftinstr{Roughly 1,100 characters. Standard eligibility runs in one
direction: the protocol tests the participant and returns eligible or screen
failure. Under the author's proposed legislation the participant is also doing
the selecting, and the second direction leaves a record. Cite
\texttt{hr9510billv5} and \texttt{paibillprior}. Adapt the criteria from
\texttt{inputs/phase-1-trial-protocol.zip},
\texttt{sections/sec-05-population.tex} \S5.1 and \S5.2.}

\begin{pafig}[mermaid-type: flowchart TD, two decision columns]
\node[mmgoal] (a) at (0,0) {Figure 14 slot\\two-way eligibility gate};
\node[mmin] (b) at (0,-2.0) {Source:\\mermaid/fig-14-eligibility-self-selection.md};
\draw[mmedge] (a) -- (b);
\end{pafig}
\vspace{-0.7cm}
\figcaption{Figure 14. Eligibility drawn as two decision columns running against each
other, what the protocol asks of the participant and what the participant may ask of\\
the protocol, converging on a single enrollment node that requires both to be satisfied.}

\draftinstr{Replace the Figure 14 slot with the mermaid-type two-column
flowchart from \texttt{mermaid/fig-14-eligibility-self-selection.md}. Two
\texttt{mmlane} boxes at $x = -3.4$ and $x = 4.6$; four \texttt{mmdec} diamonds
per lane at 2.4 cm pitch; the two edges out of \texttt{START} at looseness 0.75
so neither clips a lane title; the two "all yes" edges into enrollment at
looseness 0.8.}

\subsection{Inclusion and exclusion, in plain terms}

\draftinstr{Populate both tables from
\texttt{inputs/phase-1-trial-protocol.zip},
\texttt{sections/sec-05-population.tex} \S5.1 and \S5.2. Keep the second column
in the participant's own terms: an inclusion criterion is a statement about
them, not about the study.}

\vspace{0.30em}

\noindent\begin{tabularx}{\textwidth}{L{5.4cm} Y}
\toprule
\textbf{You may join if} & \textbf{Why this criterion exists} \\
\midrule
You are 18 or older & The study has no paediatric safety data and does not enrol children \\
Your tumour is confirmed PDAC on tissue & The drug and the operation are both specific to this disease \\
Your tumour carries a KRAS G12 mutation & Daraxonrasib acts on RAS(ON); without the mutation there is no rationale \\
Your tumour is resectable or borderline resectable & The operation must be able to achieve a clear margin \\
Your performance status is ECOG 0 or 1 & A Whipple is a major operation and the study will not expose a frailer participant to it \\
Your organ function is adequate & Both the drug and the operation load the liver, kidneys, and marrow \\
\bottomrule
\end{tabularx}

\vspace{0.30em}

\noindent\begin{tabularx}{\textwidth}{L{5.4cm} Y}
\toprule
\textbf{You may not join if} & \textbf{Why this criterion exists} \\
\midrule
Your disease is metastatic & A curative-intent resection is not the right operation \\
You have had prior pancreatic resection & The anatomy the platform is validated against is not present \\
You have an uncontrolled intercurrent illness & The perioperative risk cannot be attributed cleanly \\
You are pregnant or breastfeeding & The drug has no safety data in pregnancy \\
You cannot attend the follow-up schedule & The safety endpoints depend on the day-30 and day-90 visits \\
\bottomrule
\end{tabularx}

\vspace{0.30em}

\draftinstr{Add the remaining criteria from \S5.1 and \S5.2 of the parent
protocol so both tables are complete, and verify each table lands exactly on the
body measure with the 5.4cm plus Y split.}

\subsection{Was this tested on people like me?}

\draftinstr{Roughly 1,000 characters answering ChatGPT concern 12 honestly. The
honest answer for a first-in-human study of at most eighteen participants is
that the trial population is too small to establish subgroup performance, and
that what the protocol can commit to is reporting the composition. Cite
\texttt{FDADiver2017} for the age, race, and ethnicity reporting expectation and
\texttt{FDARealW2025} for real-world performance monitoring after the study.}

\vspace{0.30em}

\noindent\begin{tabularx}{\textwidth}{L{4.7cm} L{3.4cm} Y}
\toprule
\textbf{Characteristic} & \textbf{Reported} & \textbf{What the study can and cannot say} \\
\midrule
Age, sex, race, ethnicity & Yes, per participant & Composition is reported; subgroup performance is not estimable at $n \leq 18$ \\
Body mass index & Yes, per participant & Relevant to port placement and to force at the tip \\
Prior therapy & Yes, per participant & Prior radiation changes tissue planes and is recorded \\
Tumour location and vessel contact & Yes, per participant & Determines how much of the no-fly corridor is in play \\
Comorbidity burden & Yes, per participant & Feeds the Clavien-Dindo interpretation at day 30 \\
\bottomrule
\end{tabularx}

\vspace{0.30em}

\subsection{Booking a slot yourself, and getting an answer from a person}

\draftinstr{Roughly 1,000 characters introducing Figure 15. Under current
practice a participant learns a trial exists from a clinician who happens to
know about it. H. R. 9510 v5 proposes a direct channel with a stated service
level. State the target: five business days from first request to a confirmed
screening visit, with a 72-hour provisional hold so nobody is rushed. State the
protection: every branch notifies the treating oncologist, because the channel
is additional and not a replacement. Cite \texttt{hr9510billv5} and
\texttt{paiautospons}.}

\begin{pafig}[plantuml-type: sequence, @startuml]
\node[umlkey] (a) at (0,0) {Figure 15 slot\\booking and sponsor response};
\node[umlbox] (b) at (0,-1.9) {Source:\\plantuml/fig-15-booking-sponsor-response.puml};
\draw[umlarrow] (a) -- (b);
\end{pafig}
\vspace{-0.7cm}
\figcaption{Figure 15. Booking a screening slot directly and receiving an answer from a
named human at the sponsor, with activation bars showing who is blocked at each step\\
and every branch of the eligibility fragment notifying the participant's own oncologist.}

\draftinstr{Replace the Figure 15 slot with the PlantUML-type sequence from
\texttt{plantuml/fig-15-booking-sponsor-response.puml}. Five lifelines at
$x = 0, 3.8, 7.6, 11.4, 15.2$; \texttt{umlact} activation bars spanning exactly
the messages they enclose; message labels above the arrow, never on it; the
single self-message as a 4 mm loop at looseness 6; the \texttt{alt} fragment as
a 0.5pt \texttt{protogray} rectangle with a corner tab and a dashed divider.}

\subsection{The five choices you keep}

\draftinstr{Roughly 800 characters introducing Figure 16. A consent form
presents every choice at once at the moment of greatest stress. The same five
choices, revealed as layers that accumulate rather than replace, are each live
at a different point in the study, and not one of them is exercised on the day
of surgery. Cite \texttt{NCIConsent24} and \texttt{cfr050paiuni}.}

\begin{pafig}[d2-type: layers and steps]
\node[d2key] (a) at (0,0) {Figure 16 slot\\five choices, layered};
\node[d2box] (b) at (0,-1.8) {Source:\\d2/fig-16-five-choices-layers.d2};
\draw[d2edge] (a) -- (b);
\end{pafig}
\vspace{-0.7cm}
\figcaption{Figure 16. The five choices the participant keeps, revealed as accumulating
layers with the choice, the interval in which it is live, and what happens if the answer\\
is no, so that no decision that matters is taken at the moment of least capacity.}

\draftinstr{Replace the Figure 16 slot with the D2-type layered figure from
\texttt{d2/fig-16-five-choices-layers.d2}. Five \texttt{d2step} panels stacked
at 3.6 cm pitch; earlier panels redrawn as \texttt{d2ghost} 0.4pt outlines with
no text so the reveal reads as accumulation; the panel-to-panel connector on the
left margin at looseness 0.7.}
